% Generated draft: requires mathematical and source review.
\documentclass[11pt]{article}
\usepackage[margin=25mm]{geometry}
\usepackage{fontspec}
\setmainfont{Latin Modern Roman}
\usepackage{amsmath,amssymb,mathtools,unicode-math}
\setmathfont{Latin Modern Math}
\usepackage{enumitem,needspace,xurl,hyperref}
\setlength{\parindent}{0pt}
\setlength{\parskip}{5pt}
\setlength{\emergencystretch}{3em}
\newcommand{\N}{\mathbb{N}}
\newcommand{\Z}{\mathbb{Z}}
\newcommand{\Q}{\mathbb{Q}}
\newcommand{\R}{\mathbb{R}}
\newcommand{\C}{\mathbb{C}}
\newcommand{\F}{\mathbb{F}}
\newcommand{\A}{\mathbb{A}}
\newcommand{\PP}{\mathbb{P}}
\newcommand{\E}{\mathbb{E}}
\newcommand{\eps}{\varepsilon}
\newcommand{\dd}{\,\mathrm{d}}
\begin{document}
\section*{UnsolvedMath ID: 20001917}
\textbf{Source problem code:} catalogue entry
\subsection{Definitions and mathematical statement}
Schmidt's Problem 8.6.1 says ``Generalize this statement to higher dimensions.'' The preceding Gaidukov statement is absent from the supplied record, so its exact hypotheses and asymptotic conclusion cannot be reconstructed uniquely. The text says Gaidukov's proof is intrinsically two-dimensional and suggests Mather's minimizing measures, citing [Mather1991] and [Con-Itu1999], with the support of a minimal measure replacing a single closed geodesic.

The paired aid reports a strong lifted conclusion for closed strictly negatively curved manifolds in every dimension: asymptotic geodesics converge exponentially after Busemann phase matching. The exact original Gaidukov formulation is still missing. For general metrics it identifies an obstruction to replacing free homotopy data by classical Mather rotation data. For the kinetic Lagrangian
\[
 L(x,v)=\tfrac12\|v\|_g^2,
\]
rotation vectors of invariant probability measures lie in \(H_1(M;\mathbb R)\). At rotation zero, the minimum average kinetic action is zero and the union of supports of minimizing measures is the zero section of \(TM\). A nontrivial null-homologous free homotopy class consequently loses its conjugacy-class information. Attraction merely to the projection of that zero section is vacuous, while attraction in phase space by a unit-speed geodesic to the zero section is impossible.

The aid therefore proposes retaining a conjugacy-class-sensitive minimizing-axis set \(\mathcal A_{[\gamma]}\), and asking whether its stable basin projects onto every base point. Schematically the proposed question is
\[
 \pi\bigl(W^s(\mathcal A_{[\gamma]})\bigr)=M.
\]
The precise axis-minimization notion, treatment of deck transformations, and whether stable attraction means convergence to a set or to a phase-matched individual axis require definition. This is a proposed reduction, not a completed general higher-dimensional theorem.

The record also concatenates Section 8.7, ``Systolic and diastolic inequalities for surfaces,'' communicated by Guth, Rotman and Sabourau. For a compact Riemannian surface it uses \(\operatorname{sys}(M,g)\) for the shortest nontrivial closed geodesic. Since it applies this to \(S^2\), ``nontrivial'' here must distinguish nonconstant geodesics rather than require a noncontractible loop; the latter systole would not give its sphere quantity.

For Morse functions \(f:M\to\mathbb R\), the two source minimax expressions are
\[
 L_a(M,g)=\inf_f\max_{t\in\mathbb R}
                 \mathcal F_a(f^{-1}(t)),\qquad
 L_b(M,g)=\inf_f\max_{t\in\mathbb R}
                 \mathcal F_b(f^{-1}(t)),
\]
where \(\mathcal F_a\) is total level-set length and \(\mathcal F_b\) is the length of its longest component. The source reports realization by unions of closed geodesics and calls \(L_a\) one version of the diastole, while warning that different notions exist [Bal-Sab2010]. Critical-level conventions and the exact realization theorem are not supplied; replacing a maximum by a supremum may be necessary without an attainment statement.

The source attributes to [Croke1988], with improvements [Nab-Rot2002], [Sabourau2004] and [Rotman2006], the sphere inequality
\[
 \operatorname{sys}(S^2,g)\leq\sqrt{32}\,
                          \sqrt{\operatorname{Area}(S^2,g)}.
\]
It reports a conjectured smaller constant, printed in the extraction as ``\(\sqrt2\sqrt3\)'', and says two congruent equilateral triangles glued along their boundaries obstruct any smaller value. That radical string may represent \(\sqrt{2\sqrt3}\) or \(\sqrt2\,\sqrt3\); the nesting is not resolved by this source text and is not silently chosen here. The doubled-triangle example is nonsmooth and motivates metrics with conical singularities.

For that enlarged class the source considers
\[
 \sigma(g)=\frac{\operatorname{Area}(S^2,g)}
                 {\operatorname{sys}(S^2,g)^2}.
\]
It says \(\sigma\) is Lipschitz for an appropriate metric-space distance, asserts a positive minimum attributed to [Croke1988], conjectures that the doubled equilateral triangle attains it, and calls the round sphere a critical point in connection with the large family of Zoll metrics having the same value [Balacheff2006]. The metric distance, admissible singularities, attainment/compactness and critical-point notion are not supplied. A positive systolic lower ratio by itself gives a positive infimum, not a proof of attainment. These adjacent surface questions and contextual claims are distinct from Schmidt's higher-dimensional asymptotic-geodesic problem.
\subsection{Short English statement}
The main problem asks for an unspecified Gaidukov generalization; the aid explains why ordinary homological Mather measures lose essential free-homotopy information. It proposes a conjugacy-sensitive stable-basin formulation. A concatenated surface-systolic section is retained in full scope, including its minimax definitions, references, doubled-triangle conjecture and unresolved extracted constant.
\subsection{Sources}
\begin{itemize}
\item Local source record: catalogue entry, catalogue entry; canonical \texttt{aim-geometry-notes.json}, \texttt{notes[254]}, source item 8.6, Problem 8.6.1 attributed to Schmidt. Title: ``Higher-dimensional asymptotic geodesics and the homology obstruction.''
\item Keith Burns and Vladimir Matveev, \emph{Open problems and questions about geodesics}, AIM workshop ``Geodesics,'' supplied reference \url{https://aimath.org/pastworkshops/geodesicsproblems.pdf}. This draft uses the supplied record without accessing a PDF. The trailing page number and author header are extraction artifacts.
\item Main-problem references: Gaidukov's theorem/proof, [Mather1991] and [Con-Itu1999]. The preceding theorem is missing, and no full entries are supplied. Research attempt 1 gives the negative-curvature special case, zero-rotation Mather obstruction and proposed axis-set reduction.
\item Concatenated Section 8.7 is communicated by Guth, Rotman and Sabourau and cites [Bal-Sab2010], [Croke1988], [Nab-Rot2002], [Sabourau2004], [Rotman2006] and [Balacheff2006]. All labels are retained; full bibliography is absent.
\item Registry dated 14 August 2026: ``partially\_solved,'' ``PARTIAL-PROGRESS,'' statement status ``exact''; attempt partial, research-run status open, rigor ``proved,'' novelty confidence low. The machine-generated, non-peer-reviewed aid expressly requires independent verification before citation.
\end{itemize}
\end{document}
